%-------------------------------------------------------------------------------
%  Spoken script for the seminar "Chiral soliton lattice in inhomogeneous
%  magnetic fields".  Companion to csl_seminar.tex.
%
%  Compile:  pdflatex csl_script  (twice)
%-------------------------------------------------------------------------------
\documentclass[11pt,a4paper]{article}

\usepackage{lmodern}
\usepackage[T1]{fontenc}
\usepackage[margin=2.4cm,top=2.2cm,bottom=2.4cm]{geometry}
\usepackage{amsmath,amssymb,bm}
\usepackage{xcolor}
\usepackage{microtype}
\usepackage[colorlinks=true,linkcolor=myblue,urlcolor=myblue]{hyperref}
\usepackage{enumitem}
\usepackage{titlesec}
\usepackage{fancyhdr}

\definecolor{myblue}{HTML}{1F4E79}
\definecolor{myred}{HTML}{C0392B}
\definecolor{mygray}{HTML}{5A5A5A}
\definecolor{myorange}{HTML}{E08214}

\newcommand{\B}{\bm{B}}
\newcommand{\Hb}{\bm{H}}
\newcommand{\gr}{\bm{\nabla}}
\newcommand{\fpi}{f_\pi}
\newcommand{\mpi}{m_\pi}
\newcommand{\Bcsl}{B_{\mathrm{CSL}}}

% --- one slide entry ----------------------------------------------------------
\newcommand{\slide}[3]{%
  \par\vspace{4.5mm}%
  \noindent{\color{myblue}\rule{\textwidth}{0.9pt}}\par\vspace{-2.0mm}%
  \noindent{\color{myblue}\bfseries\large #1}%
  \hfill{\color{myorange}\bfseries #2}\par\vspace{0.4mm}%
  \noindent{\color{mygray}\small\itshape #3}\par\vspace{1.6mm}%
}
\newcommand{\cue}[1]{\par\vspace{1.0mm}\noindent{\color{mygray}\small$\blacktriangleright$\ \itshape #1}\par\vspace{1.0mm}}
\newcommand{\cut}{{\color{myred}\bfseries[CUTTABLE]}\ }
\newcommand{\act}[1]{\par\vspace{7mm}\noindent
  {\color{myred}\bfseries\Large #1}\par\vspace{-1mm}
  \noindent{\color{myred}\rule{\textwidth}{1.6pt}}\par\vspace{2mm}}

\setlength{\parindent}{0pt}
\setlength{\parskip}{2.2mm}

\pagestyle{fancy}\fancyhf{}
\fancyhead[L]{\small\color{mygray} CSL in inhomogeneous magnetic fields --- spoken script}
\fancyhead[R]{\small\color{mygray}\thepage}
\renewcommand{\headrulewidth}{0.3pt}

\begin{document}

\begin{center}
  {\LARGE\bfseries\color{myblue} Chiral soliton lattice in inhomogeneous magnetic fields}\\[2mm]
  {\large Spoken script for a 30--40 minute theory seminar}\\[3mm]
  {\color{mygray}Companion to \texttt{csl\_seminar.pdf} --- 25 content slides, 6 backup.
   Target: \textbf{39 minutes} as written, \textbf{31} with the three cuttable slides dropped.}
\end{center}

\vspace{2mm}
\hrule
\vspace{3mm}

{\color{mygray}\small
\textbf{How to use this.} The text in roman type is what to say --- not to be read
verbatim, but it is written in speaking register, so reading it aloud once or twice
will fix the phrasing in your head. The grey italic lines marked $\blacktriangleright$
are stage directions: what to point at, what to write, what to skip. The orange
times are cumulative clock, so you can tell at a glance whether you are ahead or
behind. Three slides are marked \cut: drop them in that order if you are running
long, and the talk still works.

\textbf{Pace.} Roughly 140 words a minute. The two places to deliberately slow down
are \emph{The answer: $\varphi_0$ is the magnetic scalar potential} and
\emph{Two theorems} --- that pair is the core of the new work, and it is worth
letting the room chew on it. The two places to speed up are Act I and Act II:
this audience knows ChPT, and you are reminding, not deriving.}

%===============================================================================
\act{Opening}

\slide{Title slide}{0:00}{Wait for the room to settle before starting.}

Thanks for having me. I want to talk about what the QCD ground state does in a
very strong magnetic field --- and in particular, what it does when that magnetic
field is \emph{not} uniform, which is the situation in every physical system where
the effect could conceivably matter.

This is work with Tom\'a\v{s} Brauner in Stavanger; it is on the arXiv as
2608.27319.

The shape of the talk is this. I am going to spend about ten minutes reminding you
of a story you have probably heard --- the chiral soliton lattice --- because the
new results are, quite literally, theorems about that story. Then I will spend
twenty minutes on what we actually did.

\cue{Advance.}

%-------------------------------------------------------------------------------
\slide{The setting: QCD in a field of $10^{19}$ gauss}{0:45}{1 min 15 s}

Strong-field QCD stopped being exotic some time ago. We have magnetar interiors,
with fields somewhere between $10^{15}$ and $10^{19}$ gauss; we have off-central
heavy-ion collisions, which produce the strongest magnetic fields ever made in a
laboratory; and we have lattice QCD, where a background magnetic field is now a
routine tool.

The naive expectation is that you need a field of order $m_\rho^2/e$, about
$10^{20}$ gauss, before QCD notices. That is the scale at which the chiral
condensate itself starts to be affected.

The point I want to start from is that the \emph{anomaly} gives you a much lower
scale. This one:

\cue{Point at $\Bcsl = 16\pi \fpi^2\mpi/\mu$.}

For a baryon chemical potential of order a GeV this is about $0.06$ GeV squared,
which is $10^{19}$ gauss --- an order of magnitude below the naive estimate. And
notice that it is proportional to the pion mass, so it \emph{vanishes in the chiral
limit}. That is the real surprise: in the chiral limit, an arbitrarily weak
magnetic field changes the structure of dense matter.

Above that field the ground state of QCD at finite baryon density is neither the
vacuum nor nuclear matter. It is a spatially periodic, parity-breaking condensate
of neutral pions --- the chiral soliton lattice.

And here is the fine print, which is the reason I am giving this talk. Every word
of that result assumes the magnetic field is uniform. In every system I just
listed, it is not.

%-------------------------------------------------------------------------------
\slide{The story in four steps}{2:00}{45 s --- do not linger}

So, four steps.

First, a \emph{wall}: Son and Stephanov showed in 2007 that the anomaly gives a
neutral-pion domain wall a baryon number proportional to $B$, which makes it
competitive with nuclear matter.

Second, a \emph{lattice}: Brauner and Yamamoto solved the effective theory exactly
in 2017 and turned the wall into a soliton lattice with an analytic phase diagram.

Third --- and this is the talk --- a \emph{bent field}. What survives when the
magnetic field is inhomogeneous?

And fourth, what we learned, and what we still cannot do.

If you only remember one question from the next forty minutes, make it this one:
does the chiral soliton lattice need a \emph{uniform} magnetic field, or only
\emph{some} magnetic field?

%===============================================================================
\act{Act I --- The anomaly builds a wall\hfill\normalsize\mdseries (target: 8:00 at the end)}

\slide{Why only $\pi^0$ survives}{2:45}{1 min 45 s --- this is a reminder, move briskly}

Let me set up the effective theory. This is leading-order two-flavour chiral
perturbation theory --- the usual thing, with $\Sigma$ the $SU(2)$ matrix of pions
coupled to an external electromagnetic field through the covariant derivative.

The magnetic field does two things to this theory.

The first is kinematic: it reorganises the charged pions into Landau levels, so
their energy picks up $(2n+1)eB$. Once $\sqrt{eB}$ is well above the pion mass they
are gapped and they decouple from the low-energy physics. That leaves a single real
degree of freedom --- the neutral pion phase --- and the theory collapses to
sine-Gordon.

The second thing is the interesting one. The magnetic field couples to that
remaining field only through the Wess--Zumino--Witten term. It is the same term
that gives you $\pi^0 \to 2\gamma$; here it shows up as this last piece of the
Hamiltonian.

\cue{Point at $-\Hb\cdot\gr\varphi$ and at the definition $\Hb = \mu\B/4\pi^2\fpi^2$.}

Everything in the rest of this talk follows from this one Hamiltonian density.
Gradient squared, a periodic potential, and an anomalous term linear in the
gradient. And I will use $\Hb$ throughout for the rescaled field --- it absorbs the
chemical potential and the pion decay constant, and it has dimensions of a momentum.

One word on validity, since somebody will ask. This generates a new momentum scale,
$p_{\rm CSL} = B/4\pi^2\fpi^2$, and the derivative expansion needs that to be small
compared to $4\pi \fpi$. Everything I show you is accurate up to relative
corrections of order that ratio squared.

%-------------------------------------------------------------------------------
\slide{The anomalous term is a total derivative --- and that is the point}{4:30}{2 min 30 s}

Now look at the structure of that last term. Since the magnetic field is
divergence-free, $-\Hb\cdot\gr\varphi$ is a total derivative. So it does
\emph{not} appear in the equation of motion at all.

But it does appear in the \emph{energy}. And --- this is the crucial bit --- it is
\emph{linear} in the gradient, whereas the kinetic term is quadratic. So for small
gradients the linear term always wins: it always pays to wind the pion field along
the magnetic field. The only thing stopping you is the mass term, which wants
$\varphi$ to sit at zero.

The anomaly also tells you what that winding \emph{means} physically. The anomalous
baryon current gives you a baryon density proportional to $\B\cdot\gr\varphi$, and
a magnetisation proportional to $\mu\,\gr\varphi$. So winding the pion field is a
way of storing baryon number.

\cue{Move to the right-hand column.}

Son and Stephanov put those two facts together. Take a single domain wall --- the
field goes from zero to $2\pi$ over a thickness of order the pion Compton
wavelength. Its surface energy density is $8\fpi^2\mpi$, and by the anomaly it
carries a baryon number per unit area of $eB/2\pi$.

Divide one by the other and you get the energy per baryon stored in the wall:
$16\pi \fpi^2\mpi/eB$. That falls with the magnetic field. So once the field is
strong enough, storing your baryons in pion walls is cheaper than storing them in
nucleons --- and that happens above $B_1$, about $1.1\times10^{19}$ gauss.

There is a second scale, $B_0 = 3\mpi^2/e$, which is where the wall becomes
\emph{locally} stable. Below it there is a tachyonic charged-pion mode living on
the wall; the magnetic field lifts it into a Landau level and stabilises it. That
number is $1.0 \times 10^{19}$ gauss --- essentially the same as $B_1$ for the
physical pion mass.

And again: both scales are proportional to powers of the pion mass. Both vanish in
the chiral limit.

%-------------------------------------------------------------------------------
\slide{From one wall to a stack}{7:00}{1 min}

Here is the obvious next thought. If one wall is cheaper per baryon than nuclear
matter, then so are two walls, and so are $N$ walls. So the ground state ought to
be a \emph{stack} of parallel walls, with the spacing set by the competition
between the mass term, which wants $\varphi = 0$, and the anomalous term, which
wants winding.

That is a perfectly good energetics argument, and it is where Son and Stephanov
left it. But it is a conjecture, not a solution. Two questions were open: what
exactly \emph{is} the ground state, and where exactly does it sit in the
$\mu$--$B$ plane?

That is step two.

%===============================================================================
\act{Act II --- The lattice, solved exactly\hfill\normalsize\mdseries (target: 15:00 at the end)}

\slide{The pendulum in disguise}{8:00}{2 min}

Since the magnetic field picks out a direction, the ground state can only be
modulated along it, and the equation of motion is the static sine-Gordon equation.

\cue{Say the next sentence slowly --- it is the punchline of the slide.}

Which is the equation of motion of a simple pendulum, with $z$ playing the role of
time.

That is a solved problem. The solution is the Jacobi amplitude function, with one
free parameter --- the elliptic modulus $k$ --- and the period of the lattice is
$2kK(k)$ over the pion mass.

\cue{Point at the figure, then at the curves one at a time.}

What is plotted here is the \emph{gradient} of the pion field, which by the anomaly
is the baryon density. So this picture is literally a picture of where the baryons
sit.

As $k$ goes to one --- the red curve --- you get isolated, widely separated spikes:
individual domain walls, far apart. As $k$ decreases you get a smooth, almost
sinusoidal winding. So the single Son--Stephanov wall is one limit of a
one-parameter family, and the lattice interpolates between ``dilute walls'' and
``smooth winding''.

Note that $k$ is not fixed by the equation of motion. Every $k$ solves it. The
ground state is selected by minimising the \emph{energy}, which is the next slide.

%-------------------------------------------------------------------------------
\slide{Selecting the ground state}{10:00}{1 min 45 s}

Plug the solution back in, average over one period, and you get this expression for
the energy density --- complete elliptic integrals of the first and second kind,
and a term linear in the magnetic field.

Minimise with respect to $k$ and essentially everything cancels. You are left with
one remarkably simple condition:

\cue{Point at the boxed gap equation. Pause.}

$E(k)/k = B/\Bcsl$.

And here is the whole existence proof, in one line. The function $E(k)/k$ is
bounded below by one, for $k$ between zero and one. So this equation has a solution
if and only if $B$ is bigger than $\Bcsl$. Below that field, no lattice; above it,
a lattice. And $\Bcsl$ is the scale I quoted on the first slide.

\cue{Point at the two panels.}

The left panel is $k$ in the ground state. It goes to one right at the transition
and falls off afterwards. The right panel is the period. It \emph{diverges} at the
critical field --- so just above $\Bcsl$ you have infinitely dilute, well-separated
domain walls, and as you crank up the field they pack together and the profile
becomes sinusoidal.

%-------------------------------------------------------------------------------
\slide{What makes it a genuine phase}{11:45}{2 min}

Why is this a phase and not just a field configuration?

Two reasons. First, the charges per unit cell --- baryon number and magnetic moment
--- depend only on the boundary values of $\varphi$, not on the profile. They are
topological.

Second, the state breaks symmetries: parity, and continuous translations along the
field. Breaking translations gives you a phonon with a linear dispersion relation,
and that is a genuine Nambu--Goldstone mode of this state. It is also a nice example
of hierarchical symmetry breaking --- chiral symmetry breaks first and gives you the
pions, and then the pion dynamics breaks translations and gives you the phonon, with
a parametrically small scale in between.

\cue{Move to the right column.}

And the phase has an upper edge as well. If you look at charged-pion fluctuations in
the lattice background, the equation you get is a Lam\'e equation with $n=2$, whose
spectrum is three bands. The bottom of the lowest band is this expression, and it
turns \emph{negative} at large magnetic field --- because at large $B$ the elliptic
modulus falls off like $1/B$, so the second term grows like $B^2$.

So at strong enough field the charged pions Bose condense and the neutral-pion
lattice is no longer the right description. In the chiral limit that happens at
$16\pi^4 \fpi^4/\mu^2$.

%-------------------------------------------------------------------------------
\slide{The phase diagram we inherited}{13:45}{1 min 15 s}

Putting both conditions together gives you this phase diagram. I have recomputed
these curves rather than copied them, so you can read the axes: chemical potential
in MeV, magnetic field in GeV squared on the left, gauss on the right.

The dashed line is $\Bcsl$ and the solid line is the charged-pion BEC boundary; the
shaded region in between is where the chiral soliton lattice is the ground state.
Notice where it sits: right in the magnetar range.

What I want to emphasise is how good this result is. It is leading-order chiral
perturbation theory plus the anomaly, so it is model-independent, with a controlled
error estimate --- and it is completely analytic.

\cue{Pause. Then, deliberately:}

And every line on this plot assumes a uniform magnetic field in infinite volume.
Which is wrong in every application. Heavy-ion fields vary over the transverse size
of the fireball and die in a fraction of a fermi over $c$. Magnetar fields are
dipolar. And the lattice is a finite torus.

That is what we set out to fix.

%===============================================================================
\act{Act III --- What happens when the field bends\hfill\normalsize\mdseries (target: 37:00 at the end)}

\slide{Three reasons to care about non-uniform $\B$}{15:00}{1 min 30 s}

Three reasons, in increasing order of concreteness.

Heavy-ion collisions: the field varies on the scale of the fireball and is violently
time-dependent. Compact stars: the geometry is dipolar; the variation is slow on QCD
scales, but the \emph{geometry} matters, because the field lines close on themselves,
and I will come back to that.

The third reason is the one I would emphasise for this audience.

\cue{Point at item 3.}

On a torus the magnetic flux is quantised, so you cannot tune $B$ continuously in a
lattice simulation. The standard workaround --- this is Levkova and DeTar --- is to
make the field constant on half the lattice and \emph{opposite} on the other half.
The total flux is then zero and you can dial $B$ to whatever you like.

But that means the field being simulated is a \emph{magnetic domain wall}. So the
question ``what does the CSL do in a domain-wall magnetic field?'' is not an
academic one; it is the question of what a lattice calculation would actually see.

Our questions, then: does inhomogeneity destroy the lattice? If not, for which
fields does it survive? And could a cleverly shaped field ever \emph{help}?

Same effective theory, same classical minimisation --- but now in a finite domain
$V$, with a general field $\Hb(\bm{x})$.

%-------------------------------------------------------------------------------
\slide{First: a finite box already changes things}{16:30}{2 min}

Before touching inhomogeneity, we have to deal with finite volume, because for a
non-uniform field the infinite-volume limit may not even make sense.

So: minimise the same energy functional on an interval, and --- this is the
important methodological point --- \emph{impose no boundary condition by hand}.

If you do that, the variational principle hands you its own boundary condition:

\cue{Point at the box. Say it twice if necessary.}

The normal derivative of $\varphi$ equals the normal component of $\Hb$ on the
boundary. This is the natural, Neumann-type boundary condition, and it is the single
most important technical ingredient in everything that follows.

Three consequences. First, the ground state is \emph{never} uniform in a finite box,
for any non-zero field, because the solution has to bend to meet this condition at
the walls. Second, instead of a continuum of allowed $k$ --- which is what you have
in infinite volume --- you now get a \emph{discrete} set, because only certain $k$
satisfy the boundary condition. Third, the solutions come in two families: gradient
maximal at the centre, or minimal at the centre.

\cue{Point at the figure.}

Here are both, for a box of five pion Compton wavelengths. The dotted lines are the
value $\bar H$, and you can see every curve hitting them exactly at the edges ---
that is the natural boundary condition. The blue curve is a weak field: the profile
is nearly flat in the middle and only bends near the walls. The red curve is a
strong field: there is now a genuine soliton sitting in the middle of the box.

%-------------------------------------------------------------------------------
\slide{The finite-volume phase boundary}{18:30}{1 min 30 s}

That gives us a sensible definition of the transition in a finite box: it is where
the ground state flips from the first family to the second --- from bending only at
the boundary, to carrying a soliton in the bulk. It is the finite-volume analogue of
the domain wall that signals the onset of the phase in infinite volume.

\cue{Point at the figure.}

This is the result: critical field against box size, both dimensionless. Two things
to read off it.

One, at large volume it converges to $4/\pi$, about $1.27$ --- which is exactly
$\Bcsl$ in these units. Good; the infinite-volume answer is recovered.

Two, at small volume the critical field \emph{rises}, steeply. A small box suppresses
the lattice, which makes sense: you need a stronger field to fit a soliton into a
box that is only a couple of Compton wavelengths across.

This matters if you ever want to compare with a lattice simulation, where the volume
is finite by construction. I should say that finite-volume effects in this problem
were first looked at in Tiyasa Kar's master's thesis.

%-------------------------------------------------------------------------------
\slide{The chiral limit: the problem becomes linear}{20:00}{1 min 30 s}

Now let us turn on inhomogeneity. Start with the easy case: set the pion mass to
zero.

The mass term is what made the problem non-linear, so without it the equation of
motion is just the Laplace equation --- with the same Neumann boundary condition we
just derived.

\cue{Pause.}

And that is a classic problem. The Neumann problem for the Laplace equation has a
solution, and it is \emph{unique up to an additive constant}.

Think about what that means. In the chiral limit there is exactly one stationary
point of the energy functional, so it is automatically the absolute minimum. No
competing solutions, no first-order transition, no comparing energies of branches.
The entire problem reduces to a question about a linear map: given a magnetic field,
what is the pion profile it produces?

%-------------------------------------------------------------------------------
\slide{The answer: $\varphi_0$ is the magnetic scalar potential}{21:30}{2 min 30 s}

Here is the answer, and it is clean.

Decompose the magnetic field on the domain into a gradient piece and a solenoidal
piece --- this is the Helmholtz-type decomposition adapted to a finite domain, so
$\bm{A}$ is divergence-free \emph{inside} $V$ and tangential \emph{on} the boundary.

Now complete the square in the energy.

\cue{Walk through the algebra on the slide, out loud.}

You add and subtract $\gr\psi$, and you are left with a perfect square, minus half
the gradient of $\psi$ squared, minus a cross term with $\bm A$. And the cross term
integrates to zero --- precisely because $\bm A$ has no divergence inside and no
normal component on the boundary. Integrate by parts and both surface and bulk
pieces vanish.

So the energy is bounded below by minus a half $\int(\gr\psi)^2$, and then, in a
second step, by minus a half $\int \Hb^2$.

Two saturation conditions, and they are the whole physics of this talk.

The first bound is saturated if and only if the gradient of $\varphi$ equals the
gradient of $\psi$. So the ground state \emph{is} the magnetic scalar potential. The
map from magnetic fields to pion profiles that I asked about on the last slide is
just: take the scalar potential.

The second bound is saturated if and only if $\bm A$ vanishes identically --- that
is, if and only if the magnetic field is \emph{conservative}.

%-------------------------------------------------------------------------------
\slide{Two theorems, and a picture}{24:00}{2 min}

Let me state those as two theorems and then show you what they mean.

\cue{Point at block 1.}

First: when is there \emph{no} lattice? The energy is always less than or equal to
zero, and equality happens exactly when $\gr\psi$ vanishes --- which happens exactly
when the magnetic field is tangential to the boundary \emph{everywhere}. That is
panel (c): closed field lines that never leave the region. For any other field, the
trivial vacuum is unstable and you get pion condensation. Any field that pokes
through the boundary condenses pions.

\cue{Point at block 2.}

Second: when is the lattice \emph{optimal}? When the field is conservative --- panel
(a). Then the gradient of the pion field equals the magnetic field \emph{pointwise},
the energy saturates its absolute lower bound, and, importantly, the answer does not
depend on the shape of the domain at all. The minimisation is done locally,
point by point.

\cue{Look up at the audience for this sentence.}

So the headline is: inhomogeneity by itself is not the enemy. \emph{Closed field
lines} are. A field can bend as much as it likes, and as long as it leaves the
region, the pions condense.

And for the intermediate case, panel (b), which is neither conservative nor
tangential --- there, both the decomposition and the ground state depend sensitively
on the domain. Which is to say: for realistic fields, boundaries are physical.

%-------------------------------------------------------------------------------
\slide{\cut A worked family: separation of variables}{26:00}{1 min --- do not read the formula}

Quickly, for completeness. If you take a two-dimensional field pointing everywhere
along $z$ with a magnitude that depends on $x$, on a rectangle, then finding the
scalar potential is an exercise in separation of variables --- it is electrostatics
on a rectangle --- and you get an explicit Fourier series.

For the magnetic domain wall, where $H$ is a sign function, only the odd harmonics
survive and the answer collapses to this. And if you let the box get wide, the sum
becomes an integral.

The reason I show you this is not the formula. It is that these are exact,
closed-form answers in a non-trivial inhomogeneous field --- which means we have
benchmarks to test the numerics against, and I will use them in a few slides.

%-------------------------------------------------------------------------------
\slide{\cut Bonus: a non-uniform chemical potential comes for free}{27:00}{1 min}

One generalisation that costs nothing. Suppose the chemical potential is also a
function of position --- which is what you would want for a system in local
equilibrium.

Then the rescaled field $\Hb$ is no longer divergence-free, and the Laplace equation
becomes a Poisson equation, with a source that is $\B\cdot\gr\mu$.

But look back at the derivation: nowhere did we use that $\Hb$ was solenoidal. The
decomposition still holds, $\bm A$ is still divergence-free and tangential, and both
bounds go through verbatim. So every statement survives, provided you phrase it in
terms of $\Hb$ rather than $\B$.

In particular --- and this one is worth saying --- the lattice still beats the
vacuum unless $\B$ is tangential to the boundary, and that conclusion is completely
independent of how the chemical potential varies.

%-------------------------------------------------------------------------------
\slide{Massive pions: back to non-linearity, but two facts survive}{28:00}{2 min}

Now the physically interesting case: pion mass back on, inhomogeneous field. The
equation of motion is non-linear again and in general we will have to go numerical.
But two things can still be said exactly.

\cue{Point at Fact 1.}

The first is a no-go. The mass term contributes non-negatively to the energy, so the
energy of any configuration is bounded below by its value in the chiral limit. And
we just showed that for tangential fields \emph{that} is non-negative. So tangential
fields never produce a lattice --- at any pion mass. The no-go survives.

\cue{Point at Fact 2.}

The second is a guarantee, and it is a variational argument. Use the scalar
potential $\psi$ as a trial state. Its energy is this: minus a half the gradient
squared, plus the mass term. For a conservative field the gradient of $\psi$ is just
$\Hb$, so the integrand is negative everywhere as soon as $|\Hb|$ exceeds twice the
pion mass.

Translate that into the physical magnetic field and you get $|\B| > \frac{\pi}{2}
\Bcsl$, about $0.09$ GeV squared.

\cue{Emphasise.}

Two things about that. First, it is remarkably close to the uniform-field threshold
--- a factor of $\pi/2$. Second, it holds for \emph{any} domain, whatever its shape
or size, because for conservative fields the trial state does not know about the
boundary.

I should be careful about the logic: this is a variational bound, so it is a
\emph{sufficient} condition for the lattice, not a phase boundary. That is why it
can sit above the curve I showed you earlier without contradicting it --- that curve
compared two actual solutions; this compares a trial state against the vacuum.

%-------------------------------------------------------------------------------
\slide{Numerics: what we do, and why you should believe it}{30:00}{2 min}

For everything else we minimise the energy functional directly. Let me tell you how,
because the ``why should I believe it'' question is fair.

\cue{Run down the left column quickly.}

We discretise the domain on a grid; central differences in the bulk, one-sided at
the edges; and we integrate with the trapezoidal rule --- deliberately, because
higher-order integrators develop artefacts on modulated states. Then we minimise
over all the grid values at once.

And again, no boundary condition is imposed. The natural boundary condition has to
come out of the free minimisation, which is a strong check on the whole setup.

We cross-check a global optimiser on a coarse grid against a local one on a fine
grid.

\cue{Point at the figure.}

Here is the one-dimensional benchmark. Points are the free minimisation, the line
underneath is the analytic solution. They lie on top of each other. And look at the
edges --- the gradient comes out to exactly $\bar H = 2$ there. That is the natural
boundary condition, found rather than imposed.

To check that this is not a lucky parameter choice, we scanned box size and field
strength on a fifty-by-fifty grid and asked whether the code found the true global
minimum. It did, at 2495 of the 2500 points. The five failures are all on curves
where the functional has degenerate minima, which is exactly where you would expect
a global optimiser to have trouble.

%-------------------------------------------------------------------------------
\slide{Two dimensions: a Gaussian blob of field}{32:00}{1 min 15 s}

Now two dimensions, and a genuinely non-trivial field: a localised Gaussian blob,
generated from a Gaussian scalar potential, in the chiral limit.

What is plotted is the anomalous baryon density. Bright is positive baryon density,
dark is negative.

\cue{Trace the four lobes with the pointer.}

And this is the first picture of what a soliton lattice looks like in a bent field.
It is not a stack of flat planes any more. The condensate textures itself to the
field --- it is a curved object, defined locally by the geometry of $\B$.

This is an honest minimisation on a $51\times51$ grid with no ansatz; we did not
assume a profile and fit parameters.

%-------------------------------------------------------------------------------
\slide{The case study: a magnetic domain wall}{33:15}{1 min 30 s}

Now the case I flagged at the beginning --- the lattice-QCD field. Constant, and
opposite, on the two halves, smoothed over a width $\xi$ with a hyperbolic tangent.
Tuning $\xi$ interpolates between a sharp wall and a field that is essentially
linear across the whole box.

\cue{Point at the figure.}

The headline is on the left: the lattice \emph{survives}. And note that this field
is \emph{not} conservative --- it is not in the nice class from the theorem --- and
the condensate forms anyway.

You can see the baryon density vanishing along the wall at $x = 0$, where the field
changes sign, and the pattern is symmetric about it.

\cue{Point at the benchmark block.}

And this is where the analytic work earns its keep. For a sharp wall on a square
domain in the chiral limit, the analytics say the average energy density should be
exactly half the uniform-field value. The code gives $0.499$, on a coarse grid, and
independently of the box size. That is a genuine two-dimensional check with a
non-trivial number in it.

%-------------------------------------------------------------------------------
\slide{\cut How much does bending cost?}{34:45}{1 min}

Two scans, quickly.

On the left, energy against pion mass for several aspect ratios. The intercepts at
zero pion mass reproduce the analytic values --- the dashed lines --- which is
another check. And increasing the pion mass monotonically \emph{reduces} the
condensation energy, as it must: the mass term fights the winding.

On the right, energy against box size for several wall widths. If the box is smaller
than the wall width, the field is weak everywhere inside and nothing condenses. Once
the box is much bigger than the wall width, you saturate at the uniform-field answer.

So: bending costs you condensation energy, but it never costs you all of it.

%-------------------------------------------------------------------------------
\slide{The analytic gem: how much energy lives \emph{on} the wall?}{35:45}{1 min 30 s}

I want to finish the physics with my favourite result in the paper, because it
answers ``how much does the wall cost'' exactly rather than numerically.

In the chiral limit you can do the sum. The average energy density, normalised to
the uniform-field value, is this sum over odd harmonics of a $\tanh$ --- and it
depends only on the \emph{aspect ratio} of the box, not on its size.

Now, to get anything out of that you use this summation identity:

\cue{Point at it.}

which some of you will recognise as a fermionic Matsubara sum. It turns up here in
what is otherwise a magnetostatics problem.

And you get a closed form: one, minus $14\alpha\zeta(3)/\pi^3$, plus corrections
that are exponentially small. Those first two terms are the \emph{entire} power
series in the aspect ratio --- everything else is non-analytic.

\cue{Point at the plot.}

Two limits to check. As the aspect ratio goes to zero --- a thin slab --- you
recover the uniform-field answer exactly. And for a square domain you get exactly
one half, which is the number the numerics reproduced two slides ago.

And here is the physics in that formula: the deficit is \emph{linear} in $\alpha$,
which means the total energy excess scales like $L_z^2$ with no factor of $L_x$. It
is a \emph{surface} term. The cost of the magnetic domain wall is localised on the
domain wall --- which is what you would want, but now we know it exactly.

%===============================================================================
\act{Act IV --- Where this leaves us\hfill\normalsize\mdseries (target: 39:00 at the end)}

\slide{Summary}{37:15}{1 min 15 s}

Three things to take away.

One: inhomogeneity does not kill the chiral soliton lattice. \emph{Geometry}
decides. If the field is tangential to the boundary everywhere, there is no
condensation at any pion mass. For anything else, the vacuum is unstable.

Two: conservative fields are the friendly case. In the chiral limit the pion
gradient just follows the field lines, the energy saturates its bound, and the
answer does not care about the domain. With massive pions we can still \emph{guarantee}
a lattice above about $0.09$ GeV squared, for any domain.

Three: bending costs, but not much, and it is a surface effect. For the domain-wall
field that lattice simulations actually use, the cost is exactly computable in the
chiral limit and localised on the wall.

And along the way we mapped the finite-volume phase boundary, and built a
two-dimensional minimiser that imposes no boundary conditions by hand.

%-------------------------------------------------------------------------------
\slide{What we cannot do yet}{38:30}{1 min}

Four honest gaps.

The excitation spectrum. The ground state is one property; transport and finite
temperature need the fluctuations. In a uniform field that is the Lam\'e problem.
In a bent field the phonon is not a plane wave, and it is not even obvious to me
what the Goldstone counting should be.

Realistic fields. Ours are mathematically clean. A dipole, or an event-by-event
heavy-ion field, is the honest next step, and that is what would turn this into
phenomenology.

Fluctuations. Thermal and quantum corrections already move the phase boundary in the
uniform case; nobody has done that with inhomogeneity.

And a different line of attack: this Hamiltonian discretises very cleanly onto a
qubit register, which would let you ask about real-time \emph{formation} of the
lattice rather than just its ground state. There is a backup slide if anyone wants
to talk about it.

%-------------------------------------------------------------------------------
\slide{Thank you}{39:30}{}

Thank you.

\cue{Leave the standout slide up. It is the one-sentence summary, and it is a good
thing for the room to be looking at during questions:}

Bending the magnetic field costs condensation energy, but only closing the field
lines destroys the lattice.

%===============================================================================
\newpage
\act{Questions you should expect}

{\color{mygray}\small Ordered roughly by how likely they are from a QCD audience.
Answers are short on purpose --- give the one-line version first and let them push.}

\textbf{1. Why are you allowed to drop the charged pions?}\\
Landau quantisation: their gap goes like $\sqrt{eB}$, so there is a window
$\sqrt{eB} \gg \mpi$ where they are heavy and the neutral pion is the only light
mode. And I am not claiming they never matter --- they come back, and that is
exactly what sets the upper edge of the phase, $B_{\rm BEC}$. Backup slide has the
Lam\'e calculation.

\textbf{2. The WZW term also has a piece with a time derivative. Doesn't that
change the charged-pion dispersion?}\\
It does change the dispersion, but not the critical field, because $B_{\rm BEC}$ is
extracted from $\omega = 0$ where that term drops out. This is worked out in
appendix C of Brauner--Yamamoto.

\textbf{3. Is the classical treatment justified? What about fluctuations?}\\
The loop expansion is controlled by $p_{\rm CSL}^2/(4\pi\fpi)^2$, which is small in
the regime of interest. Thermal and quantum corrections have been computed for the
uniform case and they do shift the boundary --- Brauner, Kole\v{s}ov\'a and
Yamamoto, and the NLO paper. For inhomogeneous fields, nobody has done it, and I
would not want to guess.

\textbf{4. $10^{19}$ gauss is far above observed magnetar surface fields. So is this
relevant to anything?}\\
Honest answer: not yet, and we say so in the paper. Interior fields can be much
larger than surface fields, but I would not hang an argument on that. The concrete
near-term application is the lattice one --- the domain-wall field is a real
simulation setup, and now there is a prediction for what it should show.

\textbf{5. Is your two-dimensional domain-wall field actually divergence-free?}\\
Yes. It points along $z$ and depends only on $x$, so the divergence vanishes
identically. It is not \emph{conservative} --- its curl is non-zero --- which is
exactly why it is an interesting test case.

\textbf{6. A hard boundary at QCD length scales is unphysical. How seriously should
I take the boundary conditions?}\\
Seriously, and that is part of the message. The domain stands for the region where
the field is supported, or for the simulation volume. Our result is precisely that
for non-conservative fields the answer \emph{does} depend on that choice, whereas for
conservative fields it does not. So the sensitivity is a physical statement, not an
artefact --- and it tells you which questions are well posed.

\textbf{7. How do you know your minimiser found the global minimum and not a local
one?}\\
Three ways. The one-dimensional scan against the exact solution: 2495 of 2500 points,
with the failures on degeneracy curves. The two-dimensional benchmark against the
analytic domain-wall result: $0.499$ against $0.5$. And a global optimiser on a
coarse grid cross-checking the local optimiser on a fine one.

\textbf{8. What about the continuum limit?}\\
The theory has one intrinsic scale, the pion Compton wavelength, so the criterion is
that the cell be smaller than that: $n \gtrsim \mpi L$ in each direction. A concrete
example is on a backup slide --- a 1\% error on a $31 \times 31$ grid for parameters
where the criterion is comfortably satisfied. A true continuum extrapolation is
expensive because the number of optimisation variables grows as $n_x n_z$.

\textbf{9. Where does the baryon number of a \emph{finite} wall come from? The
anomalous term is a total derivative.}\\
That is Son and Stephanov's analysis. For a finite wall the total-derivative piece
integrates to zero and the baryon number comes from the Skyrme term instead: the
configuration is topologically a Skyrmion with $N_B = eBR^2/2$, and the edge of the
wall is a superconducting string of charged pions.

\textbf{10. How does this compete with the other candidate ground states --- charged
pion condensates, domain-wall Skyrmions, crystalline baryonic matter?}\\
In a uniform field there is a whole literature on exactly that, and several of those
states win in parts of the phase diagram. We have not put any of them in an
inhomogeneous field, and I think that is a genuinely open and interesting problem.

\textbf{11. Isn't your lower bound just Cauchy--Schwarz?}\\
The bound itself is elementary --- it is completing the square. The content is in the
\emph{saturation conditions}: one of them identifies the ground state as the magnetic
scalar potential, and the other singles out conservative fields. The inequality is
the easy half.

\end{document}
